\section{词义与多义性问题}

\subsection{一词多义是常态}

自然语言中的绝大多数词都有\textbf{多个含义}。Manning 以 ``pike'' 为例：

\begin{figure}[H]
\centering
\includegraphics[width=0.95\textwidth]{slides-images/slide-030.jpg}
\caption{pike 的多种含义：尖头武器、鱼类、收费公路、跳水姿势等\protect\footnotemark}
\end{figure}
\footnotetext{Slides 第31页。}

\begin{itemize}
    \item 一种\textbf{武器}（长矛/戟）
    \item 一种\textbf{鱼类}（梭鱼）
    \item 一种\textbf{道路}（收费公路，turnpike 的缩写）
    \item 一种\textbf{跳水姿势}
    \item 动词：用矛刺
    \item 澳大利亚俚语：临阵退缩
\end{itemize}

更常见的多义词如 bank（银行 / 河岸）、star（星星 / 明星）等在日常文本中无处不在。

\begin{warningbox}{词义消歧的困难}
虽然字典会为每个词列出明确的义项（sense 1, sense 2, ...），但实际上词义之间的边界往往是\textbf{模糊的}。例如 ``field'' 可以指农田、运动场、学术领域、数学中的域等——这些含义之间有程度不等的关联，硬性划分义项本身就是一种``人为的简化''。
\end{warningbox}

\subsection{Word2Vec 如何处理多义词？}

Word2Vec 为每个词只学习\textbf{一个}向量，不区分不同义项。那么这个唯一的向量代表什么？

\begin{importantbox}{词义的叠加（Superposition）}
一个多义词的词向量是其各个义项向量的\textbf{加权平均}，权重为各义项在语料中出现的相对频率：
$$v_{\text{pike}} = f_1 \cdot v_{\text{pike}_1} + f_2 \cdot v_{\text{pike}_2} + \cdots + f_k \cdot v_{\text{pike}_k}$$
其中 $f_i$ 是第 $i$ 个义项的相对频率，$v_{\text{pike}_i}$ 是该义项如果单独训练会得到的向量。
\end{importantbox}

这种``叠加''看起来会丢失信息——从一个平均向量怎么可能恢复各个义项？但有一个惊人的数学结果：

\begin{knowledgebox}{稀疏编码与义项恢复}
由于词向量空间是\textbf{高维且稀疏}的，来自\textbf{稀疏编码理论}的方法可以从单个叠加向量中\textbf{恢复出各个义项的向量}。Arora 等人（其中包括 Stanford 的 Tengyu Ma）在论文中展示了这一技术：从 ``tie'' 的单一向量中成功分离出了``领带''、``比赛平局''、``电线束扎''、``音乐连接符''等不同义项。
\end{knowledgebox}

不过在实际应用中，现代 NLP 更倾向于使用\textbf{上下文词向量}（contextual word vectors，如 BERT、GPT），它们为每个词在每次出现时生成不同的向量，从根本上解决了多义性问题。

\subsection{本章小结}

一词多义是自然语言的基本特征。Word2Vec 等静态词向量将多义词表示为各义项的加权平均，虽然看似丢失了信息，但高维空间的稀疏性使得义项在理论上可以被恢复。这一发现为理解词向量的表示能力提供了深刻的数学洞察。后续课程中将介绍的上下文词向量（如 BERT）则从根本上解决了多义性问题。

%% ========== Part 6: 分类与神经网络入门 ==========

